\section{Implementation}

\subsection{Signal}

\subsubsection{Die Klasse}

Ein digitales Signal wird in seiner simpelsten Form als eine sich ständig
ändernde Variable dargestellt. Der Datentyp der Variable wird dem Nutzer 
überlassen, indem die Signalklasse als Template implementiert wird. 

\begin{lstlisting}[language=C++, 
                       frame=single, 
                       caption={Signalklasse Version 0.1},
                       captionpos=b]
template<typename T>
class Signal {
  T value;
}
\end{lstlisting}

Da viele Funktionen, vor allem 
Filter, auf vorherige Werte zu greifen, lohnt es sich diese auch zu speichen.
Dafür bietet sich ein RingBuffer an. Ein RingBuffer ist ein fest definierter
Speicherblock (zum Beispiel ein Array) mit fester Größe, das ist vorallen 
auf Mikrocontrollern sehr 
wichtig, da hier der Speicher sehr begrenzt ist und Speicherallokation während 
der Laufzeit schnell zu Abstürzen führen kann. Ein Ringbuffer hat auserdem 
einen Index Wert welcher anzeigt, wo in diesen Speicherblock geschrieben 
werden soll. Der Index wird bei jedem schreiben erhöht und fängt wieder am Anfang
des Speicherblocks an zu schreiben, wenn das Ende erreicht ist.
Die Größe des Speichers soll auch den Nutzer überlassen werden. Sie muss aber zur
Compilezeit festgelegt sein.

\begin{lstlisting}[language=C++, 
                       frame=single, 
                       caption={Signalklasse Version 0.2},
                       captionpos=b]
template <typename T, unsigned int BUFFER_SIZE = 256>
class Signal {
private:
  T buffer[BUFFER_SIZE] = {0};
  unsigned int index = 0;
}
\end{lstlisting}

Wenn das Signal zu einem neuen Wert übergeht, muss der Index erhöht werden.
Das passiert in einer Inputfunktion, welche mit dem Prefix \glqq from-\grqq\ 
gekennzeichnet sind. Eine Inputfunktion gibt immer eine Referenz auf das 
Signalobjekt zurück.

\begin{lstlisting}[language=C++, 
                       frame=single, 
                       caption={Signalklasse Version 0.3},
                       captionpos=b]
template <typename T, unsigned int BUFFER_SIZE = 256>
class Signal {
private:
  T buffer[BUFFER_SIZE] = {0};
  unsigned int index = 0;
public:
  Signal& fromValue(T value) {
    index = (index + 1) % BUFFER_SIZE;
    buffer[index] = value;
    return *this;
  }
}
\end{lstlisting}

Eine Funktion, die welche mit dem Prefix \glqq return-\grqq\ 
gekennzeichnet sind, ist eine Ausgangsfunktion. 
Solche Funktionen geben keine Signalreferenz zurück, sondern zum Beispiel den 
momentanen Wert.

\begin{lstlisting}[language=C++, 
                       frame=single, 
                       caption={Signalklasse Version 0.3},
                       captionpos=b]
template <typename T, unsigned int BUFFER_SIZE = 256>
class Signal {
private:
  T buffer[BUFFER_SIZE] = {0};
  unsigned int index = 0;
public:
  ...

  T returnValue() {
    return buffer[index];
  }
}
\end{lstlisting}

Alle anderen Funktionen sind Funktionen, welche eine Signalreferenz zurückgeben
und im gegansatz zu Inputfunktionen nicht den Index erhöhen. Das sind zum Beispiel
add, multiply oder applyFunction. Letzteres erlaubt es dem Nutzer selbstdefinierte 
in die Verarbeitungskette einzufügen.
\begin{lstlisting}[language=C++, 
                       frame=single, 
                       caption={Signalklasse Version 0.4},
                       captionpos=b]
template <typename T, unsigned int BUFFER_SIZE = 256>
class Signal {
private:
  T buffer[BUFFER_SIZE] = {0};
  unsigned int index = 0;
public:
  ...

  Signal& add(Signal &other) {
    this->buffer[index] += other.buffer[other.index];
    return *this;
  }

  Signal& multiply(T value) {
    this->buffer[index] = T(this->buffer[index] * value);
    return *this;
  }

  SIgnal applyFunction(T (*func)(T data, void *param), 
                       void *param) {

    buffer[index] = func(buffer[index], param);
    return *this;
  }
}
\end{lstlisting}

\subsubsection{Anwendung}

\subsubsection{Anmerkungen}
Dies ist auch ein erster, großer unterschied zu

\subsection{Sound}

 I2S funktioniert mit 
Pulse-Code-Modulation (PCM) Daten. Das ist ein unkomprimiertes Audio format mit
meistens 16,24 oder 32 Bit unsigned Integern.  In dem config Header ist ein 
Datentyp mit dem Namen sample definiert, welcher genau so ein unsigned Int ist.
Hier kann man die Auflösung global einstellen.
% todo finde Quellen zu PCM.

